problem:  
Let \( M^6 \) be a compact, simply connected Riemannian manifold of **nonnegative sectional curvature** with an **effective isometric \( S^1 \)-action**.  Let \( Q = M^6 / S^1 \) be the resulting **5‑dimensional Alexandrov space** of curvature ≥ 0, whose singular set \( \Sigma \subset Q \) is the image of nonprincipal orbits.  Assume that \( Q \) is homeomorphic to \( S^5 \).  

It is known that components of the fixed‑point set \( M^{S^1} \) are even‑dimensional, and that strata arising from finite isotropy correspond to **codimension‑2 or higher** extremal subsets in the Alexandrov sense.  Suppose that such a stratum \( \Sigma^{3} \subset Q \) is homeomorphic to the 3‑sphere \( S^3 \) and locally modeled as an extremal subset of codimension 2; i.e. we assume the quotient admits local neighborhoods where \( Q \) is metrically a cone over a spherical suspension of a circle action’s link.  

**Question (Conditional Curvature–Rigidity Conjecture for Codimension 2 Extremal Subsets):**  
If a 5‑dimensional Alexandrov space \( Q \cong S^5 \) of curvature ≥ 0 arises as the quotient of a simply connected, nonnegatively curved 6‑manifold by an isometric \( S^1 \)-action, and \( Q \) contains a codimension‑2 extremal subset \( \Sigma^3 \) homeomorphic to \( S^3 \),  
must \( \Sigma^3 \) be **unknotted** in \( S^5 \)?  

That is, is it possible that \( \Sigma^3 \) represents a nontrivial smooth embedding \( S^3 \hookrightarrow S^5 \) (a genuinely high‑dimensional “knot”) while still being realizable as a curvature‑bounded extremal stratum compatible with the quotient geometry of a nonnegatively curved \( 6 \)-manifold with isometric circle symmetry?  

Equivalently, does curvature ≥ 0 and the orbit‑space structure prohibit any Alexandrov‑compatible metric near a knotted codimension‑2 stratum in \( S^5 \)?  

Why is it a “good” problem:  
This formulation resolves previous structural ambiguities by **explicitly treating the codimension‑2 stratum homeomorphic to \( S^3 \) as a conditional hypothesis**, rather than assuming a universally fixed isotropy model.  It thereby defines a sharp, well‑posed question that remains open and compelling:

* **Mathematical soundness:** It no longer presupposes a particular slice representation; instead it asks conditionally about the possible topology of a codimension‑2 extremal subset with clearly stated topological type.  
* **Conceptual simplicity:** The problem asks a succinct geometric question—“can a nonnegatively curved Alexandrov \( S^5 \) arising from an isometric circle quotient contain a knotted \( S^3 \) stratum?”—that is conceptually graspable yet nontrivial.  
* **Research depth and cross‑field connection:** The question bridges  
  - **Alexandrov geometry** (extremal subset structure and curvature comparison),  
  - **Riemannian transformation group theory** (orbit‑space and isotropy analysis), and  
  - **high‑dimensional knot theory** (embeddings \( S^3 \hookrightarrow S^5 \) and their isotopy classes).  
  Its resolution would bring together techniques from each of these areas.  
* **Potential impact:** A positive answer would suggest a powerful curvature‑driven rigidity phenomenon—nonnegative curvature excludes topologically knotted orbit strata—while a counterexample would reveal new geometric models of Alexandrov spaces exhibiting hidden flexibility.  
* **Elegant character:** Despite its technical depth, the question is clean and visual: it inquires whether curvature restrictions alone can enforce *topological triviality of a singular stratum*—a vivid instance of profound geometry expressed in simple form.